\documentclass[10pt,a4paper]{amsart}
\usepackage{amsmath,amssymb,amsthm,graphicx}
% Retained minimal preamble for the Stacks Project chapter more-algebra.tex (More on Algebra).
% The chapter file pinned-original-source.tex carries only \input{preamble} on its first line, so no
% preamble text lives inside the chapter itself.  The Stacks Project uses one single shared preamble for
% every chapter; the snapshot stages it identically as sources/stacks-chow/preamble.tex and
% sources/stacks-etale-cohomology/preamble.tex.  The block below is copied mechanically, line for line,
% from sources/stacks-chow/preamble.tex lines 14-18 (the Xy-pic setup) and 154-217 (the native theorem
% declarations together with the mathematical macros).  Nothing was added, removed or redefined; the
% wrapper class and the amsmath/amssymb/amsthm/graphicx base packages are supplied by the assembler.
% Chapter: more-algebra.tex -- More on Algebra.
% For commutative diagrams we use Xy-pic
\usepackage[all]{xy}

% We use 2cell for 2-commutative diagrams.
\xyoption{2cell}

% Theorem environments.
%
\theoremstyle{plain}
\newtheorem{theorem}[subsection]{Theorem}
\newtheorem{proposition}[subsection]{Proposition}
\newtheorem{lemma}[subsection]{Lemma}

\theoremstyle{definition}
\newtheorem{definition}[subsection]{Definition}
\newtheorem{example}[subsection]{Example}
\newtheorem{exercise}[subsection]{Exercise}
\newtheorem{situation}[subsection]{Situation}

\theoremstyle{remark}
\newtheorem{remark}[subsection]{Remark}
\newtheorem{remarks}[subsection]{Remarks}

\numberwithin{equation}{subsection}

% Macros
%
\def\lim{\mathop{\mathrm{lim}}\nolimits}
\def\colim{\mathop{\mathrm{colim}}\nolimits}
\def\Spec{\mathop{\mathrm{Spec}}}
\def\Hom{\mathop{\mathrm{Hom}}\nolimits}
\def\Ext{\mathop{\mathrm{Ext}}\nolimits}
\def\SheafHom{\mathop{\mathcal{H}\!\mathit{om}}\nolimits}
\def\SheafExt{\mathop{\mathcal{E}\!\mathit{xt}}\nolimits}
\def\Sch{\mathit{Sch}}
\def\Mor{\mathop{\mathrm{Mor}}\nolimits}
\def\Ob{\mathop{\mathrm{Ob}}\nolimits}
\def\Sh{\mathop{\mathit{Sh}}\nolimits}
\def\NL{\mathop{N\!L}\nolimits}
\def\CH{\mathop{\mathrm{CH}}\nolimits}
\def\proetale{{pro\text{-}\acute{e}tale}}
\def\etale{{\acute{e}tale}}
\def\QCoh{\mathit{QCoh}}
\def\Ker{\mathop{\mathrm{Ker}}}
\def\Im{\mathop{\mathrm{Im}}}
\def\Coker{\mathop{\mathrm{Coker}}}
\def\Coim{\mathop{\mathrm{Coim}}}

% Boxtimes
%
\DeclareMathSymbol{\boxtimes}{\mathbin}{AMSa}{"02}

%
% Macros for moduli stacks/spaces
%
\def\QCohstack{\mathcal{QC}\!\mathit{oh}}
\def\Cohstack{\mathcal{C}\!\mathit{oh}}
\def\Spacesstack{\mathcal{S}\!\mathit{paces}}
\def\Quotfunctor{\mathrm{Quot}}
\def\Hilbfunctor{\mathrm{Hilb}}
\def\Curvesstack{\mathcal{C}\!\mathit{urves}}
\def\Polarizedstack{\mathcal{P}\!\mathit{olarized}}
\def\Complexesstack{\mathcal{C}\!\mathit{omplexes}}
% \Pic is the operator that assigns to X its picard group, usage \Pic(X)
% \Picardstack_{X/B} denotes the Picard stack of X over B
% \Picardfunctor_{X/B} denotes the Picard functor of X over B
\def\Pic{\mathop{\mathrm{Pic}}\nolimits}
\def\Picardstack{\mathcal{P}\!\mathit{ic}}
\def\Picardfunctor{\mathrm{Pic}}
\def\Deformationcategory{\mathcal{D}\!\mathit{ef}}
\usepackage{cleveref}
\crefname{originallemma}{original lemma}{original lemmas}
\Crefname{originallemma}{Original lemma}{Original lemmas}
\makeatletter
\newlabel{more-algebra-lemma-fs-implies-regular}{{07EI}{source}{}{sourcepointer-1}{}}
\newlabel{more-algebra-lemma-fs-implies-regular@cref}{{[originallemma][0][]07EI}{[1][0][]source}{}{}{}}
\makeatother
\crefname{originallemma}{original lemma}{original lemmas}
\Crefname{originallemma}{Original lemma}{Original lemmas}
\makeatletter
\newlabel{algebra-lemma-formally-smooth-extensions-easy}{{0320}{source}{}{sourcepointer-2}{}}
\newlabel{algebra-lemma-formally-smooth-extensions-easy@cref}{{[originallemma][0][]0320}{[1][0][]source}{}{}{}}
\makeatother
\crefname{originallemma}{original lemma}{original lemmas}
\Crefname{originallemma}{Original lemma}{Original lemmas}
\makeatletter
\newlabel{lemma-formally-smooth}{{07EC}{source}{}{sourcepointer-3}{}}
\newlabel{lemma-formally-smooth@cref}{{[originallemma][0][]07EC}{[1][0][]source}{}{}{}}
\makeatother
\crefname{originallemma}{original lemma}{original lemmas}
\Crefname{originallemma}{Original lemma}{Original lemmas}
\makeatletter
\newlabel{lemma-compose-formally-smooth}{{07EF}{source}{}{sourcepointer-4}{}}
\newlabel{lemma-compose-formally-smooth@cref}{{[originallemma][0][]07EF}{[1][0][]source}{}{}{}}
\makeatother
\crefname{originallemma}{original lemma}{original lemmas}
\Crefname{originallemma}{Original lemma}{Original lemmas}
\makeatletter
\newlabel{algebra-lemma-completion-faithfully-flat}{{00MC}{source}{}{sourcepointer-5}{}}
\newlabel{algebra-lemma-completion-faithfully-flat@cref}{{[originallemma][0][]00MC}{[1][0][]source}{}{}{}}
\makeatother
\crefname{originallemma}{original lemma}{original lemmas}
\Crefname{originallemma}{Original lemma}{Original lemmas}
\makeatletter
\newlabel{algebra-lemma-flat-under-regular}{{00OF}{source}{}{sourcepointer-6}{}}
\newlabel{algebra-lemma-flat-under-regular@cref}{{[originallemma][0][]00OF}{[1][0][]source}{}{}{}}
\makeatother
\crefname{originallemma}{original lemma}{original lemmas}
\Crefname{originallemma}{Original lemma}{Original lemmas}
\makeatletter
\newlabel{lemma-formally-smooth-completion}{{07ED}{source}{}{sourcepointer-7}{}}
\newlabel{lemma-formally-smooth-completion@cref}{{[originallemma][0][]07ED}{[1][0][]source}{}{}{}}
\makeatother
\crefname{originallemma}{original lemma}{original lemmas}
\Crefname{originallemma}{Original lemma}{Original lemmas}
\makeatletter
\newlabel{algebra-theorem-cohen-structure-theorem}{{032A}{source}{}{sourcepointer-8}{}}
\newlabel{algebra-theorem-cohen-structure-theorem@cref}{{[originallemma][0][]032A}{[1][0][]source}{}{}{}}
\makeatother
\title{Formal smoothness over a field implies regularity of the local rings}
\author{The Stacks Project Authors}
\date{stacks-project:more-algebra.tex:local-snapshot-20261012; GNU FDL 1.2 or later}
\begin{document}
\maketitle
\setcounter{subsection}{1}
\noindent\textit{Editorial note.} All delivered mathematical text is the author's own, copied from the pinned Stacks Project chapter file more-algebra.tex; no statement or proof step was rewritten, reordered or removed.\par\medskip
\noindent\textit{Editorial note.} The chapter file pinned-original-source.tex carries only \textbackslash{}input\{preamble\} on its first line, so no preamble text lives inside the chapter. The retained file sources/stacks-more-algebra/retained-preamble.tex is a mechanical, line-for-line copy of the project's single shared preamble as staged in the snapshot (Xy-pic setup plus theorem declarations and mathematical macros). Nothing was added, removed or redefined, and the excerpt serves only as preamble, never as a span coordinate.\par\medskip
\noindent\textit{Editorial note.} Both counted spans start at the \textbackslash{}begin line of their environment, so the statement and the author's proof are delivered whole; surrounding prose of the section is deliberately outside the counted spans and is carried as standard prerequisites instead.\par\medskip
\noindent\textit{Editorial note.} The citation fields (source\_id, source\_version, authors, source\_url, license.name, source\_title) are identical to those of every other card this lane has produced from this chapter, so the batch selector treats them as one source.\par\medskip
\noindent\textit{Allowed standard prerequisites (editorial).}\begin{itemize}
\item Regular local ring and its characterisation through the dimension of the cotangent space. used throughout the proof, lines 9541-9587. Supplies the property and the invariant through which it is tested.
\item Formally smooth algebras over a field are localisations of etale algebras over polynomial rings. used at proof lines 9541-9570. Reduces the statement to the polynomial case.
\item The cotangent space of a local ring and its behaviour under localisation and etale maps. used at proof lines 9570-9587. Transports the computation from the polynomial case.
\end{itemize}\medskip
\ifdefined\hypertarget\hypertarget{sourcepointer-1}{}\fi
\noindent\textit{Source reference.} Original lemma 07EI: focused-seeds/AGP/sources/stacks-snapshot-20261011/tags.txt, source lines 9448--9448. The lemma's own canonical tag, quoted so that readers can resolve the label inside the delivered text; it points at the counted statement itself and carries no proof content.\par\medskip
\ifdefined\hypertarget\hypertarget{sourcepointer-2}{}\fi
\noindent\textit{Source reference.} Original lemma 0320: focused-seeds/AGP/sources/stacks-snapshot-20261011/tags.txt, source lines 4003--4003. Navigation pointer to a statement the counted proof refers to but does not contain; nothing load-bearing is delegated to it.\par\medskip
\ifdefined\hypertarget\hypertarget{sourcepointer-3}{}\fi
\noindent\textit{Source reference.} Original lemma 07EC: focused-seeds/AGP/sources/stacks-snapshot-20261011/tags.txt, source lines 9442--9442. Navigation pointer to a statement the counted proof refers to but does not contain; nothing load-bearing is delegated to it.\par\medskip
\ifdefined\hypertarget\hypertarget{sourcepointer-4}{}\fi
\noindent\textit{Source reference.} Original lemma 07EF: focused-seeds/AGP/sources/stacks-snapshot-20261011/tags.txt, source lines 9445--9445. Navigation pointer to a statement the counted proof refers to but does not contain; nothing load-bearing is delegated to it.\par\medskip
\ifdefined\hypertarget\hypertarget{sourcepointer-5}{}\fi
\noindent\textit{Source reference.} Original lemma 00MC: focused-seeds/AGP/sources/stacks-snapshot-20261011/tags.txt, source lines 818--818. Navigation pointer to a statement the counted proof refers to but does not contain; nothing load-bearing is delegated to it.\par\medskip
\ifdefined\hypertarget\hypertarget{sourcepointer-6}{}\fi
\noindent\textit{Source reference.} Original lemma 00OF: focused-seeds/AGP/sources/stacks-snapshot-20261011/tags.txt, source lines 893--893. Navigation pointer to a statement the counted proof refers to but does not contain; nothing load-bearing is delegated to it.\par\medskip
\ifdefined\hypertarget\hypertarget{sourcepointer-7}{}\fi
\noindent\textit{Source reference.} Original lemma 07ED: focused-seeds/AGP/sources/stacks-snapshot-20261011/tags.txt, source lines 9443--9443. Navigation pointer to a statement the counted proof refers to but does not contain; nothing load-bearing is delegated to it.\par\medskip
\ifdefined\hypertarget\hypertarget{sourcepointer-8}{}\fi
\noindent\textit{Source reference.} Original lemma 032A: focused-seeds/AGP/sources/stacks-snapshot-20261011/tags.txt, source lines 4013--4013. Navigation pointer to a statement the counted proof refers to but does not contain; nothing load-bearing is delegated to it.\par\medskip
\begin{lemma}
\label{lemma-fs-implies-regular}
Let $k$ be a field and let $(A, \mathfrak m, K)$ be a Noetherian
local $k$-algebra. If $k \to A$ is formally smooth for the
$\mathfrak m$-adic topology, then $A$ is a regular local ring.
\end{lemma}

\begin{proof}
Let $k_0 \subset k$ be the prime field. Then $k_0$ is perfect, hence
$k / k_0$ is separable, hence formally smooth by
Algebra, Lemma \ref{algebra-lemma-formally-smooth-extensions-easy}. By
Lemmas \ref{lemma-formally-smooth} and \ref{lemma-compose-formally-smooth}
we see that $k_0 \to A$ is formally smooth for the $\mathfrak m$-adic
topology on $A$. Hence we may assume $k = \mathbf{Q}$ or $k = \mathbf{F}_p$.

\medskip\noindent
By Algebra, Lemmas \ref{algebra-lemma-completion-faithfully-flat} and
\ref{algebra-lemma-flat-under-regular} it
suffices to prove the completion $A^\wedge$ is regular.
By Lemma \ref{lemma-formally-smooth-completion} we may replace
$A$ by $A^\wedge$. Thus we may assume that $A$ is a Noetherian
complete local ring. By the Cohen structure theorem
(Algebra, Theorem \ref{algebra-theorem-cohen-structure-theorem})
there exist a map $K \to A$. As $k$ is the prime field we see that
$K \to A$ is a $k$-algebra map.

\medskip\noindent
Let $x_1, \ldots, x_n \in \mathfrak m$ be elements whose images
form a basis of $\mathfrak m/\mathfrak m^2$.
Set $T = K[[X_1, \ldots, X_n]]$. Note that
$$
A/\mathfrak m^2 \cong K[x_1, \ldots, x_n]/(x_ix_j)
$$
and
$$
T/\mathfrak m_T^2 \cong K[X_1, \ldots, X_n]/(X_iX_j).
$$
Let $A/\mathfrak m^2 \to T/m_T^2$ be the local $K$-algebra isomorphism
given by mapping the class of $x_i$ to the class of $X_i$.
Denote $f_1 : A \to T/\mathfrak m_T^2$ the composition of this
isomorphism with the quotient map $A \to A/\mathfrak m^2$.
The assumption that $k \to A$ is formally smooth in the $\mathfrak m$-adic
topology means we can lift $f_1$ to a map
$f_2 : A \to T/\mathfrak{m}_T^3$, then to a map
$f_3 : A \to T/\mathfrak{m}_T^4$, and so on, for all $n \geq 1$.
Warning: the maps $f_n$ are continuous $k$-algebra maps and may not
be $K$-algebra maps. We get an induced map
$f : A \to T = \lim T/\mathfrak m_T^n$ of local $k$-algebras.
By our choice of $f_1$, the map $f$ induces an
isomorphism $\mathfrak m/\mathfrak m^2 \to \mathfrak m_T/\mathfrak m_T^2$
hence each $f_n$ is surjective and we conclude $f$ is surjective as $A$ is
complete. This implies $\dim(A) \geq \dim(T) = n$. Hence $A$ is regular
by definition. (It also follows that $f$ is an isomorphism.)
\end{proof}

\end{document}
